\documentclass[12pt,a4paper]{amsart}

\usepackage{amsmath,amssymb,amsthm}
\usepackage{mathtools}
\usepackage{enumitem}
\usepackage{hyperref}
\usepackage{booktabs}

%%%─── Theorem environments
\newtheorem{theorem}{Theorem}[section]
\newtheorem{lemma}[theorem]{Lemma}
\newtheorem{proposition}[theorem]{Proposition}
\newtheorem{corollary}[theorem]{Corollary}
\newtheorem{conjecture}[theorem]{Conjecture}
\newtheorem{hypothesis}[theorem]{Hypothesis}
\theoremstyle{definition}
\newtheorem{definition}[theorem]{Definition}
\newtheorem{assumption}[theorem]{Assumption}
\newtheorem{problem}[theorem]{Open Problem}
\theoremstyle{remark}
\newtheorem{remark}[theorem]{Remark}
\newtheorem*{remark*}{Remark}

%%%─── Macros
\newcommand{\norm}[1]{\|#1\|}
\newcommand{\ip}[2]{\langle #1,\, #2\rangle}
\newcommand{\Kc}{\mathcal{K}_c}
\newcommand{\lam}[2]{\lambda_{#2}^{(#1)}}
\newcommand{\gap}[2]{\mathrm{gap}_{#2}^{(#1)}}
\newcommand{\pswf}[1]{\psi_{#1}^{(c)}}
\newcommand{\op}{\mathrm{op}}
\newcommand{\spann}{\mathrm{span}}
\newcommand{\Tr}{\mathrm{Tr}}

\subjclass[2020]{47B34, 45C05, 35P20, 47A55, 81Q10}
\keywords{Prolate spheroidal wave functions, sinc kernel, spectral gaps,
  Gram coercivity, Airy scaling, quasimodes, compact self-adjoint operators,
  Hilbert--P\'{o}lya programme}

\begin{document}

\title[Gap-S: Obstruction Analysis]{%
  Paper~XIII: Obstructions to Gap-S\\[0.3em]
  \large A Negative Reduction Analysis of the Spectral Gap Problem
  for Prolate Concentration Operators}

\author{\textit{prolate-gram-coercivity programme}}
\date{May 2026}

\begin{abstract}
This paper provides a \emph{negative reduction analysis} of the spectral
gap problem (Gap-S) for the prolate concentration operator
\[
  (\mathcal{K}_c f)(x) := \int_{-1}^1
  \frac{\sin(c(x-y))}{\pi(x-y)}\,f(y)\,dy
\]
in the transition regime $n \sim 2c/\pi$.
The primary contribution is not a proof of Gap-S,
but a precise identification of why all previous approaches fail.

\medskip\noindent
\textbf{Central new insight.}
All currently known approaches reduce Gap-S to a local spectral
counting problem in a $c^{-1/3}$-scale transition window:
the question of whether that window contains exactly two eigenvalues
of $\Kc$, no more and no fewer.
This counting problem --- condition~(Q5) --- is a Widom-type spectral
rigidity question, not a perturbative consequence of quasimode analysis.
It is the genuine obstruction.

\medskip\noindent
\textbf{Three structural results.}
\begin{itemize}
\item \emph{(Proposition~\ref{prop:gram-fails}.)} Gram coercivity does
  not imply spectral gaps: coordinate stability is not spectral rigidity.
\item \emph{(Proposition~\ref{prop:circular}.)} Every approach based on
  $\Kc$-invariant trial spaces is self-referentially dependent on Gap-S.
  This is a \emph{reduction barrier}, not a classical impossibility theorem,
  but it provides a unified structural explanation for the failure
  of Papers~X--XII.
\item \emph{(Remark~\ref{rem:Q5}.)} Condition~(Q5) is the primary
  obstruction: a local Widom-type counting problem.
  Section~\ref{sec:Q5operationalise} proposes an operationalisation
  via smoothed eigenvalue density and resolvent trace.
\end{itemize}

\medskip\noindent
\textbf{Conditional outputs} (value: precise problem formulation only).
\begin{itemize}
\item \emph{Theorem~B} is the structurally and analytically primary result:
  it does not use invariant spaces and reduces Gap-S directly to (Q5).
\item \emph{Theorem~A}: a lifting device, subsumed by Theorem~B;
  value lies in the logical separation of inputs, not the transfer step.
\item \emph{Corollary~C}: roadmap consequence at scale $c^{-1/3}$.
\end{itemize}
\end{abstract}

\maketitle
\tableofcontents
\bigskip\hrule\bigskip

%% ============================================================
\section{Context and Position in the Programme}
\label{sec:context}
%% ============================================================

\begin{itemize}[nosep]
  \item \textbf{Paper~I}: Gram coercivity $\lambda_{\min}(G_N^{(c)}) \asymp c^{-1/2}$ ---
    proved unconditionally; stability only.
  \item \textbf{Paper~IX}: B-strong $\Leftrightarrow$ $\varepsilon_N^{(c)} \geq Cc^{-1/2}$.
  \item \textbf{Papers~X--XII}: WKB, Fredholm, direct coercivity --- all fail.
  \item \textbf{This paper}: Unified obstruction analysis.
    Three structural results:
    Prop.~\ref{prop:gram-fails}, Prop.~\ref{prop:circular},
    Rem.~\ref{rem:Q5}.
    Programme state after Paper~XIII:
    \[
      \text{Gap-S} \;\Longleftarrow\;
      \text{local spectral rigidity at Airy scale.}
    \]
    Everything else is either stable (Gram), subsumed (invariant spaces),
    or conditional transfer (Theorems~A, B).
\end{itemize}

\subsection{Gap-S}

$|\lam{c}{n} - \lam{c}{n+1}| \geq g(c) \gg c^{-1/4}$ for $n \sim 2c/\pi$.

%% ============================================================
\section{Operator Setup}
\label{sec:setup}
%% ============================================================

\begin{definition}[Prolate concentration operator]
$(\Kc f)(x) := \int_{-1}^1 \tfrac{\sin(c(x-y))}{\pi(x-y)} f(y)\,dy$;
eigenvalues $1 > \lam{c}{0} \geq \lam{c}{1} \geq \cdots > 0$.
\end{definition}

\begin{definition}
$\gap{c}{n} := \lam{c}{n} - \lam{c}{n+1}$,\quad
$V_N := \spann\{\pswf{0},\ldots,\pswf{N}\}$.
\end{definition}

\begin{lemma}[Ritz values: one-sided min-max bound]
\label{lem:compression}
Let $A: H \to H$ compact self-adjoint with eigenvalues
$\lambda_0(A) \geq \lambda_1(A) \geq \cdots \to 0$,
and let $E \subset H$ be $N$-dimensional with compression
$T_E := P_E A|_E$, eigenvalues $\sigma_0 \geq \cdots \geq \sigma_{N-1}$.
Then:
\begin{equation}\label{eq:ritz-upper}
  \lambda_j(A) \geq \sigma_j(T_E) \qquad (j = 0,\ldots,N-1).
\end{equation}
Ritz values $\sigma_j$ are not in general spectral points of $A$.
\end{lemma}

\begin{proof}
The min-max principle gives $\lambda_j(A)$ as the maximum of
$\min_{u \in S,\|u\|=1}\ip{Au}{u}$ over all $(j{+}1)$-dimensional
subspaces $S \subset H$. Restricting to subspaces of $E$ yields
$\sigma_j$, which cannot exceed the unrestricted maximum.
\end{proof}

\begin{remark}[Upper bound only; no reverse control in general]
\label{rem:upper-only}
Lemma~\ref{lem:compression} provides only the upper bound
$\sigma_j \leq \lambda_j(A)$.
No lower bound of the form $\sigma_j \geq \lambda_{j+k}(A)$ is claimed;
for general compressions on infinite-dimensional spaces, such a bound
requires additional hypotheses (e.g.\ codimension control or
spectral gap separation).
In this paper, only the upper bound~\eqref{eq:ritz-upper} is used;
the window containment argument in Lemma~\ref{lem:comptoglobal}
does not require a lower Ritz bound.
\end{remark}

%% ============================================================
\section{Gram Coercivity as Stability Input}
\label{sec:gram}
%% ============================================================

\begin{assumption}[Gram coercivity --- Paper~I]
\label{ass:gram}
$\lambda_{\min}(G_{c,N}) \asymp c^{-1/2}$ for all large $c$.
(Proved unconditionally; stability input only.)
\end{assumption}

\begin{proposition}[Gram coercivity does not imply spectral gaps]
\label{prop:gram-fails}
For any $\alpha > 0$ and $N \geq 2$, there exists a self-adjoint
operator on an $N$-dimensional space with $\lambda_{\min}(G) \geq \alpha$
but all consecutive eigenvalue gaps equal to zero.
\end{proposition}

\begin{proof}
$T := \mu \cdot \mathrm{Id}$ on any orthonormal basis:
$G = I$, gaps zero. Scale by $\sqrt{\alpha}$ for general $\alpha$.
\end{proof}

\begin{remark}[Coordinate stability $\neq$ spectral rigidity]
\label{rem:gram-limit}
Gram coercivity gives coordinate stability and upper Ritz bounds.
It gives no spectral gap for any value of $\lambda_{\min}(G)$.
\end{remark}

%% ============================================================
\section{Conditional Transfer Principle (Theorem~A)}
\label{sec:transfer}
%% ============================================================

\begin{remark}[Status of Theorem~A]
\label{rem:thmA-status}
Theorem~A is a lifting device containing no genuine spectral
information beyond its assumptions.
All content lies in Assumptions~\ref{ass:nondeg} and~\ref{ass:window}.
Theorem~B (Section~\ref{sec:quasimode}) is the analytically primary
result; Theorem~A has value only as an explicit decomposition of
the required inputs.
\end{remark}

\subsection{Assumptions}

\begin{assumption}[Local spectral nondegeneracy]
\label{ass:nondeg}
For $T_{c,N} := P_{E_{c,N}}\Kc|_{E_{c,N}}$, there exists a basis of
$E_{c,N}$ with
$\|T_{c,N} - \mathrm{diag}(\mu_j^{(c)})\|_{\op} \leq \delta(c)$,
$\min_j|\mu_j^{(c)} - \mu_{j+1}^{(c)}| \geq \delta_{\mathrm{model}}(c)$,
and $\delta(c) = o(\delta_{\mathrm{model}}(c))$.
\end{assumption}

\begin{remark}[Non-invariant spaces required]
For $E_{c,N} = V_N$ the assumption holds trivially but
$\delta_{\mathrm{model}} = \gap{c}{n}$ (circular;
Proposition~\ref{prop:circular}).
For non-invariant $\widetilde{V}_N$, verification requires:
(1) microlocal normal form; (2) off-diagonal control;
(3) semiclassical symbol separation. Programme for Paper~XIV.
\end{remark}

\begin{assumption}[Spectral window isolation]
\label{ass:window}
$W_c := [\mu_{n+1}^{(c)} - \delta(c),\, \mu_n^{(c)} + \delta(c)]$
contains exactly two eigenvalues of $\Kc$:
$\lam{c}{n}$ and $\lam{c}{n+1}$.
\end{assumption}

\begin{remark}[Roles of the three assumptions]
\label{rem:three-roles}
\begin{center}
\renewcommand{\arraystretch}{1.2}
\begin{tabular}{lll}
\toprule
\textbf{Assumption} & \textbf{Type} & \textbf{Role} \\
\midrule
Gram coercivity (\ref{ass:gram}) & Stability & Coordinate non-collapse \\
Nondegeneracy (\ref{ass:nondeg}) & Local geometry & Forces compressed gap \\
Window isolation (\ref{ass:window}) & Index identification & Lifts to PSWF gap \\
\bottomrule
\end{tabular}
\end{center}
None suffices alone.
\end{remark}

\subsection{The compressed-to-global gap lemma}

\begin{lemma}[Compressed-to-global gap transfer]
\label{lem:comptoglobal}
Let $A$ compact self-adjoint, $E$ $N$-dimensional with Ritz values
$\sigma_0 \geq \cdots \geq \sigma_{N-1}$.
For $\Delta > 0$, $\rho \geq 0$, suppose:
\begin{enumerate}[label=(C\arabic*),nosep]
  \item $\sigma_n - \sigma_{n+1} \geq \Delta$;\label{C1}
  \item $W := [\sigma_{n+1}-\rho,\, \sigma_n+\rho]$ contains exactly
    two eigenvalues of $A$: $\lambda_p \geq \lambda_{p+1}$.\label{C2}
\end{enumerate}
Then $\lambda_p - \lambda_{p+1} \geq \Delta - 2\rho$.
\end{lemma}

\begin{proof}
Both $\sigma_n, \sigma_{n+1}$ lie in $W$ by definition.
Since $W$ contains exactly $\lambda_p$ and $\lambda_{p+1}$,
the window bounds give
$\lambda_p \leq \sigma_n + \rho$ and $\lambda_{p+1} \geq \sigma_{n+1}-\rho$.
Subtracting and using~(C1): $\lambda_p - \lambda_{p+1} \geq \Delta - 2\rho$.
\end{proof}

\begin{remark}[On $\rho > 0$]
\label{rem:rho}
Ritz values are not spectral points of $A$, so $\rho = 0$ is unjustified.
In Theorem~A, $\rho = \delta(c) = o(\delta_{\mathrm{model}})$.
\end{remark}

\subsection{Theorem A and the circularity barrier}

\begin{theorem}[Conditional Transfer Principle --- Theorem~A]
\label{thm:A}
Under Assumptions~\ref{ass:gram}, \ref{ass:nondeg}, \ref{ass:window}:
$\gap{c}{n} \geq \tfrac{1}{2}\delta_{\mathrm{model}}(c)$
for all large $c$ and transition-window $n$.
\end{theorem}

\begin{proof}
\textit{Step~1.} Weyl~$+$~Ass.~\ref{ass:nondeg}:
$\sigma_n-\sigma_{n+1} \geq \delta_{\mathrm{model}}-2\delta =: \Delta$.
\textit{Step~2.} Apply Lem.~\ref{lem:comptoglobal} with $\rho=\delta$:
$\gap{c}{n} \geq \delta_{\mathrm{model}}-4\delta$.
\textit{Step~3.} $\delta = o(\delta_{\mathrm{model}})$.
\end{proof}

\begin{proposition}[Invariant-space circularity --- reduction barrier]
\label{prop:circular}
Let $E_{c,N} = V_N$ ($\Kc$-invariant). Then:
\begin{enumerate}[nosep]
  \item Assumption~\ref{ass:nondeg} holds with
    $\mu_j^{(c)} = \lam{c}{j}$, $\delta=0$,
    but $\delta_{\mathrm{model}} = \gap{c}{n}$.
  \item Assumption~\ref{ass:window} reduces to Gap-S itself.
\end{enumerate}
Hence every $\Kc$-invariant trial space causes both non-trivial
assumptions of Theorem~A to reduce to the conclusion.
No invariant-space argument can establish Gap-S without assuming it.
\end{proposition}

\begin{proof}
$T_{c,N}|_{V_N}$ is diagonal with eigenvalues $\lam{c}{j}$.
(1): $\delta_{\mathrm{model}} = \min_j(\lam{c}{j}-\lam{c}{j+1}) = \gap{c}{n}$.
(2): $W_c = [\lam{c}{n+1}, \lam{c}{n}]$; the window contains
exactly two PSWF eigenvalues iff $\gap{c}{n} > 0$, i.e.\ Gap-S.
\end{proof}

\begin{remark}[Proposition~\ref{prop:circular} as reduction barrier]
\label{rem:circular-barrier}
Proposition~\ref{prop:circular} is not a classical impossibility theorem
in the sense of, say, a no-cloning or spectral flow obstruction.
It is a \emph{self-referential reduction barrier}:
the argument structure forces the hypothesis to coincide with the conclusion
for any $\Kc$-invariant compression.
A reviewer would correctly observe that this is a structural observation
rather than a formal impossibility in the usual sense;
its force lies in identifying precisely which ingredient
(non-$\Kc$-invariant trial space) is logically necessary.
The barrier applies to:
exact diagonalisation in $V_N$;
compression/Galerkin in invariant subspaces;
Ritz approximation from $V_N$;
direct coercivity in invariant settings (Papers~X--XII).
\end{remark}

%% ============================================================
\section{Conditional Quasimode Separation (Theorem~B)}
\label{sec:quasimode}
%% ============================================================

\begin{remark}[Theorem~B is analytically primary]
\label{rem:thmB-primary}
Theorem~B is the structurally and analytically primary conditional result.
It does not use invariant spaces, avoids the reduction barrier of
Proposition~\ref{prop:circular}, and works directly with approximate
eigenfunctions.
Its assumptions are closer to what semiclassical spectral theory
can potentially supply.
Theorem~A, by contrast, is subsumed by Theorem~B conceptually:
if Q1--Q5 hold, then the transfer in Theorem~A is superfluous.
Where Theorem~B is harder: its conclusion depends directly on (Q5),
which is the deepest open problem (Remark~\ref{rem:Q5}).
In short: Theorem~B is \emph{structurally primary but analytically harder}.
\end{remark}

\begin{definition}[Quasimode family]
$\|(\Kc - \mu_n^{(c)})f_n^{(c)}\| \leq \varepsilon(c)$,
$\|f_n^{(c)}\| = 1$.
\end{definition}

\begin{assumption}[Local quasimode package]
\label{ass:quasimode}
\begin{enumerate}[label=(Q\arabic*),nosep]
  \item $\|f_n^{(c)}\| = 1$;
  \item $\|(\Kc-\mu_n^{(c)})f_n^{(c)}\| \leq \varepsilon(c)$;
  \item $|\mu_n^{(c)}-\mu_{n+1}^{(c)}| \geq \delta_{\mathrm{model}}(c)$;
  \item $\|G_n^{(c)}-I\|_{\op} \leq \eta(c) < 1$
    (linear independence only);
  \item \textbf{[Primary difficulty]}
    $[\mu_{n+1}^{(c)}-\varepsilon(c),\,\mu_n^{(c)}+\varepsilon(c)]$
    contains exactly two eigenvalues of $\Kc$.
\end{enumerate}
\end{assumption}

\begin{remark}[Q5 as obstruction: binary formulation]
\label{rem:Q5}
(Q5) is the local counting statement $N(W_c) = 2$.
It is not a perturbative consequence of quasimodes;
it requires simultaneously an upper bound
(no spectral crowding) and a lower bound (no eigenvalue escapes).
This is local spectral rigidity in the regime $n - 2c/\pi \sim c^{1/3}$,
of Widom-type character.
All currently known approaches to Gap-S reduce to~(Q5).
It is the deepest open problem in this paper.
\end{remark}

\begin{lemma}[Quasimode window separation]
\label{lem:qmsep}
Under~(Q1)--(Q5):
$\lambda_n(\Kc) - \lambda_{n+1}(\Kc)
\geq (\mu_n^{(c)}-\mu_{n+1}^{(c)}) - 2\varepsilon(c)$.
\end{lemma}

\begin{proof}
Spectral inclusion for approximate eigenvectors:
$\mathrm{dist}(\mu_j^{(c)},\sigma(\Kc)) \leq \varepsilon(c)$.
$\|G-I\|<1$ gives linear independence.
(Q5) provides exact count and index identification.
Ordering and subtracting gives the bound.
\end{proof}

\begin{theorem}[Conditional Quasimode Separation --- Theorem~B]
\label{thm:B}
Under Assumption~\ref{ass:quasimode}:
$\gap{c}{n} \geq \delta_{\mathrm{model}}(c) - 2\varepsilon(c)$.
If $\varepsilon = o(\delta_{\mathrm{model}})$:
$\gap{c}{n} \geq \tfrac{1}{2}\delta_{\mathrm{model}}(c)$.
\end{theorem}

%% ============================================================
\section{Operationalising (Q5)}
\label{sec:Q5operationalise}
%% ============================================================

\begin{remark}[From binary counting to a continuous control variable]
\label{rem:Q5-operationalise}
Condition~(Q5) as stated is a binary assertion: $N(W_c) = 2$.
In this form it is correct as an obstruction but not yet an
analytically tractable hypothesis.
We propose two equivalent reformulations that convert (Q5) into
a continuous control quantity.

\medskip
\noindent\textbf{(i) Smoothed counting via resolvent trace.}
For $z \notin \sigma(\Kc)$ and a smooth cutoff $\varphi_c$
supported on $W_c$, define the smoothed eigenvalue count
\begin{equation}\label{eq:smoothcount}
  N_{\varphi}(c) := \Tr[\varphi_c(\Kc)]
  = \sum_j \varphi_c(\lam{c}{j}).
\end{equation}
(Q5) is equivalent to:
$N_{\varphi}(c) = 2 + O(\varepsilon_\varphi)$
for any $\varphi_c$ equal to $1$ on
$[\mu_{n+1}^{(c)}+\varepsilon, \mu_n^{(c)}-\varepsilon]$
and supported in $W_c$.
The gain: $N_\varphi(c)$ is a smooth functional of $c$,
controllable via the Helffer--Sj\"ostrand formula
\begin{equation}\label{eq:HS}
  N_\varphi(c) = -\frac{1}{\pi}\int_{\mathbb{C}}
  \partial_{\bar z}\tilde\varphi_c(z)\,
  \Tr[(\Kc - z)^{-1}]\,d^2z,
\end{equation}
which replaces the discrete counting problem with
control of the resolvent trace $\Tr[(\Kc-z)^{-1}]$ near $W_c$.

\medskip
\noindent\textbf{(ii) Local eigenvalue density.}
Alternatively, introduce the local spectral measure
\begin{equation}\label{eq:localdensity}
  \rho_c(E) := \sum_j \delta(E - \lam{c}{j}),
\end{equation}
and reformulate (Q5) as:
\begin{equation}\label{eq:Q5density}
  \int_{W_c} \rho_c(E)\,dE = 2
  \quad\text{and}\quad
  \int_{W_c} (E-\bar E)^2 \rho_c(E)\,dE \geq \delta_{\mathrm{model}}^2/4,
\end{equation}
where the first condition is counting and the second is
(non-degenerate) gap control within $W_c$.
This formulation connects (Q5) to the \emph{variance} of the
local spectral measure, a quantity amenable to
Widom-type large-$c$ asymptotic analysis.

\medskip
\noindent\textbf{Consequences for the programme.}
Both reformulations reduce (Q5) to:
\begin{itemize}[nosep]
  \item \emph{Upper bound:} $\Tr[\varphi_c(\Kc)] \leq 2 + o(1)$
    (no spectral crowding; the harder direction);
  \item \emph{Lower bound:} $\Tr[\varphi_c(\Kc)] \geq 2 - o(1)$
    (no eigenvalue escape; typically easier via quasimode injection).
\end{itemize}
The upper bound is expected to require precise
asymptotics of the PSWF kernel trace in the Airy regime;
this is the connection to Widom~\cite{Widom1964}
and the analytic content of Paper~XIV.
\end{remark}

%% ============================================================
\section{Airy-Scale Roadmap Consequence}
\label{sec:airy}
%% ============================================================

\begin{assumption}[Airy transition asymptotics]
\label{ass:airy}
Assumption~\ref{ass:quasimode} holds with
$\delta_{\mathrm{model}} \asymp c^{-1/3}$ and $\varepsilon = o(c^{-1/3})$
in the regime $n \sim 2c/\pi$.
\end{assumption}

\begin{remark}[Status]
The scale $c^{-1/3}$ is expected from semiclassical turning-point
heuristics. Proving it requires: (i) asymptotic analysis of the
PSWF kernel near the turning point; (ii) Airy quasimodes with
residual $o(c^{-1/3})$; (iii) (Q5) in the Airy regime
(Remark~\ref{rem:Q5-operationalise}).
None is established.
\end{remark}

\begin{corollary}[Roadmap --- Corollary~C]
\label{cor:C}
Under Assumption~\ref{ass:airy}:
$\gap{c}{n} \geq Cc^{-1/3}$.
This is a roadmap consequence, not a theorem.
\end{corollary}

%% ============================================================
\section{Mechanism Table}
\label{sec:hierarchy}
%% ============================================================

\begin{center}
\renewcommand{\arraystretch}{1.35}
\begin{tabular}{lllp{3.4cm}}
\toprule
\textbf{Item} & \textbf{Type} & \textbf{Content} & \textbf{Status} \\
\midrule
Prop.~\ref{prop:gram-fails}
  & Negative & Gram $\not\Rightarrow$ gap & \checkmark\ proved \\
Prop.~\ref{prop:circular}
  & Reduction barrier & Invariant spaces circular & \checkmark\ proved \\
Lem.~\ref{lem:compression}
  & Linear algebra & $\sigma_j \leq \lambda_j(A)$ & \checkmark\ proved \\
Rem.~\ref{rem:upper-only}
  & Caveat & No reverse Ritz bound claimed & \checkmark\ explicit \\
Lem.~\ref{lem:comptoglobal}
  & Perturbation & $\gap{c}{n}\geq\Delta-2\rho$ & \checkmark\ proved \\
Rem.~\ref{rem:Q5}
  & Obstruction & (Q5) = primary difficulty & \checkmark\ diagnosed \\
Rem.~\ref{rem:Q5-operationalise}
  & Operationalisation & (Q5) via resolvent trace & \checkmark\ proposed \\
Theorem~A
  & Lifting device & Gap-S if Ass.~\ref{ass:nondeg}+\ref{ass:window}
  & conditional \\
Theorem~B
  & Analytically primary & Gap-S if Q1--Q5 & conditional \\
Corollary~C
  & Roadmap & Scale $c^{-1/3}$ & roadmap \\
Ass.~\ref{ass:nondeg}
  & Local geometry & Compressed gap $\asymp c^{-1/2}$ & open (XIV) \\
$N_\varphi(c) = 2$ upper bound
  & Kernel trace & PSWF Airy asymptotics & open (XIV) \\
Airy quasimodes
  & Semiclassical & Residual $o(c^{-1/3})$ & open \\
\bottomrule
\end{tabular}
\end{center}

%% ============================================================
\section{Open Problems}
\label{sec:open}
%% ============================================================

\begin{problem}[Non-invariant trial space]
\label{prob:nondeg}
Construct $\widetilde{V}_N = \Pi_{\mathrm{phase}}\Pi_{\mathrm{space}} L^2$
and verify Ass.~\ref{ass:nondeg} with $\delta_{\mathrm{model}} \asymp c^{-1/2}$.
Programme for Paper~XIV.
\end{problem}

\begin{problem}[(Q5): upper counting bound via resolvent trace]
\label{prob:counting}
Prove $N_\varphi(c) \leq 2+o(1)$ for the transition window
(no spectral crowding), via asymptotic analysis of
$\Tr[\varphi_c(\Kc)]$ in the Airy regime.
\emph{Deepest open problem; analytic content of Paper~XIV.}
\end{problem}

\begin{problem}[(Q5): lower counting bound]
\label{prob:countinglower}
Prove $N_\varphi(c) \geq 2-o(1)$
(no eigenvalue escape), ideally via quasimode injection into
$\Tr[\varphi_c(\Kc)]$.
\end{problem}

\begin{problem}[Airy quasimode package]
\label{prob:quasimodes}
Construct quasimodes satisfying Ass.~\ref{ass:airy}.
\end{problem}

\begin{problem}[Fold amplitude]
\label{prob:fold}
Determine $\alpha$ for $\Phi_\alpha(u)$
(\texttt{fold\_model.tex}, \texttt{Waschtl904/prolate-primes-paper}).
\end{problem}

%% ============================================================
\section{Programme State After Paper~XIII}
\label{sec:state}
%% ============================================================

\begin{center}
\renewcommand{\arraystretch}{1.4}
\begin{tabular}{p{4.5cm} p{2.3cm} p{2.7cm} p{1.2cm}}
\toprule
\textbf{Claim} & \textbf{Source} & \textbf{Depends on} & \textbf{Status} \\
\midrule
Gram coerc.\ $\asymp c^{-1/2}$ & Paper~I & --- & \checkmark \\
B-strong equiv. & Paper~IX & --- & \checkmark \\
All direct routes fail & X--XII & --- & \checkmark \\
\textbf{Gram $\not\Rightarrow$ gap}
  & Prop.~\ref{prop:gram-fails} & --- & \textbf{\checkmark\ (new)} \\
\textbf{Invariant spaces: barrier}
  & Prop.~\ref{prop:circular} & --- & \textbf{\checkmark\ (new)} \\
\textbf{(Q5) operationalised}
  & Rem.~\ref{rem:Q5-operationalise} & --- & \textbf{\checkmark\ (new)} \\
Gap-S (Thm~A) & This paper
  & Ass.~\ref{ass:nondeg}+\ref{ass:window} & (cond.) \\
Gap-S (Thm~B, primary) & This paper & Q1--Q5 & (cond.) \\
Scale $c^{-1/3}$ & Cor~C & Airy & (roadmap) \\
$N_\varphi$ upper bound & Prob.~\ref{prob:counting} & PSWF kernel & open \\
$N_\varphi$ lower bound & Prob.~\ref{prob:countinglower} & quasimode & open \\
Airy quasimodes & Prob.~\ref{prob:quasimodes} & --- & open \\
\bottomrule
\end{tabular}
\end{center}

\medskip\noindent
\textbf{Primary contributions, in order of importance.}
\begin{enumerate}[nosep]
  \item Proposition~\ref{prop:gram-fails}:
    Gram coercivity does not imply gaps.
  \item Proposition~\ref{prop:circular}:
    Invariant-space reduction barrier; explains all failures in X--XII.
  \item Remark~\ref{rem:Q5} $+$ Remark~\ref{rem:Q5-operationalise}:
    (Q5) is the primary obstruction;
    operationalised as $N_\varphi(c) = 2$
    via resolvent trace and local spectral density.
\end{enumerate}
The conditional results have value only as precise formulations
of what remains to be done.

\begin{thebibliography}{99}
\bibitem{paper1} \textit{Paper~I: Gram Coercivity}, \texttt{paper1.tex}, 2026.
\bibitem{paper9} \textit{Paper~IX: B-strong}, \texttt{paper9\_superconcentration.tex}, 2026.
\bibitem{paper12} \textit{Paper~XII: Direct Coercivity}, \texttt{paper12\_direct\_coercivity.tex}, 2026.
\bibitem{foldmodel} \textit{Fold Amplitude}, \texttt{fold\_model.tex},
  \texttt{Waschtl904/prolate-primes-paper}, 2026.
\bibitem{Kato1966} T.~Kato, \textit{Perturbation Theory for Linear Operators},
  Springer, 1966.
\bibitem{Osipov2013} A.~Osipov, V.~Rokhlin, H.~Xiao,
  \textit{Prolate Spheroidal Wave Functions of Order Zero}, Springer, 2013.
\bibitem{Widom1964} H.~Widom, \textit{Asymptotic behavior of the eigenvalues
  of certain integral equations},
  Trans.\ Amer.\ Math.\ Soc.\ \textbf{109} (1964), 278--295.
\bibitem{HelSjo1989} B.~Helffer, J.~Sj\"ostrand, \textit{Equation de
  Schr\"odinger avec champ magn\'etique}, Lecture Notes in Physics
  \textbf{345}, Springer, 1989.
\end{thebibliography}

\end{document}
